\documentclass[../../main.tex]{subfiles}
\begin{document}

\begin{center}
  {\fontsize{48pt}{58pt}\selectfont \textbf{1998}\normalsize\textbf{年度}}
\end{center}

\begin{table}[htb]
  \centering
  \caption{出題テーマと難易度}
  \begin{tabular}{|c|p{9cm}|c|}
    \hline
    \textbf{問題番号} & \textbf{テーマ} & \textbf{難易度} \\
    \hline
    第1問 & 交換可能な行列と零行列 & \\
    \hline
    第2問 & 積分で定まる関数列 & \\
    \hline
    第3問 & 三角形の頂点を通る円の方程式 & \\
    \hline
    第4問 & 三角方程式の根の極限 & \\
    \hline
    第5問 & 梯子状の道の接続確率 & \\
    \hline
    第6問 & 三角関数の定積分と回転体 & \\
    \hline
  \end{tabular}
\end{table}

\begin{table}[htb]
  \centering
  \caption{解答}
  \renewcommand{\arraystretch}{1.4}
  \begin{tabular}{|c|c|p{9cm}|}
    \hline
    \textbf{問題番号} & \textbf{小問} & \textbf{答え} \\
    \hline
    第1問 & (1)，(2) &  \\
    \hline
    第2問 & -- & $f_n(x)=2^{n+1}x^2+5\cdot3^{n-1}-2^{n+1}$ \\
    \hline
    第3問 & -- & 証明問題 \\
    \hline
    第4問 & (1) & $0$ \\
    \cline{2-3}
    & (2) & $\dfrac{\sin a}{\pi}$ \\
    \hline
    第5問 & (1) & $Q_2=p(2-p)$，$Q_3=p^2(3-2p)$ \\
    \cline{2-3}
    & (2) & $Q_n=p^{n-1}\{n-(n-1)p\}$ \\
    \hline
    第6問 & (1) & $I_2=\dfrac{38-26\sqrt{2}}{105}$ \\
    \cline{2-3}
    & (2) & $V=3\pi I_3+2\pi I_2$（証明問題） \\
    \hline
  \end{tabular}
\end{table}

\paragraph{試験形態}
6問．試験時間は要確認．

\end{document}
